\clearpage\SourceLeaf{299}
\begin{center}
\textbf{[LOCUS}\par\medskip
\textbf{OF THE APEX OF A SPHERICAL TRIANGLE}\par\medskip
\textbf{ON A GIVEN BASE WITH A GIVEN AREA.]}
\end{center}

\medskip

\begin{center}
\textsc{Gauss} to \textsc{Schumacher}. Göttingen, 29 December 1841.
\end{center}

\mbox{.\,.\,.\,.\,.\,.\,.} Do you know of a \textls{printed} solution to the problem of finding the locus of the apex of a spherical triangle on a given base and with a given area? I recently had occasion to look for one and was somewhat surprised that the analytical solution leads by a somewhat roundabout route to a very simple result whose synthetic proof can be given in a few lines. \mbox{.\,.\,.\,.\,.\,.\,.}

\medskip

\begin{center}
[\textsc{Schumacher} \emph{to} \textsc{Gauss}.\ \emph{Altona, 3 January 1842.}]
\end{center}

\{A geometrical solution to the problem you mentioned, my dearest friend, was given not very long ago by \textsc{Steiner} in \textsc{Crelle}'s journal (vol.\ 2.\ p.\ 45), where he does not cite his predecessors correctly. \textsc{Lexell}'s paper is not in Nova Acta Petrop.\ V.\ Pars 1, as stated there. It is in Acta Petrop.\ pro 1781, Pars I. \textsc{Euler} returned to it later in Nova Acta Tom.\ X.\ p.\ 47 (Variae speculationes etc.). \textsc{Legendre} has it
